\section{Introduction}
\label{sec:introduction}

\IEEEPARstart{G}{lobal} freshwater security represents one of the most critical socio-ecological challenges of the twenty-first century. According to the United Nations Food and Agriculture Organization (FAO), irrigated agriculture accounts for approximately 70\% of global freshwater withdrawals, climbing beyond 85\% in arid and developing economies~\cite{allen1998crop, ahmad2024ai}. Despite this staggering resource footprint, conventional agricultural water management remains notoriously inefficient: between 40\% and 60\% of diverted irrigation water is squandered globally through uncalibrated surface flooding, fixed calendar timers, and uncompensated soil evaporation~\cite{jones2004irrigation, olatinwo2025review}. Over-irrigation not only depletes vulnerable aquifer reserves but also triggers root-zone hypoxia, fungal phytopathologies, and pervasive nitrate leaching into downstream watersheds. Conversely, under-irrigation induces acute hydraulic cavitation in xylem vessels, leaf turgor loss, stomatal closure, and irreversible yield degradation~\cite{jackson1981canopy, xu2026turgor}. Consequently, closing the loop between real-time physiological crop water demand and autonomous actuation has emerged as an imperative for climate-resilient agriculture.

Crop water dynamics are fundamentally governed by the Soil-Plant-Atmosphere Continuum (SPAC), a thermodynamic cascade wherein water potential gradients drive moisture flux from bulk soil through root membranes, xylem vascular architecture, and foliar mesophyll cells, terminating in transpirational vapor release across stomatal apertures into the turbulent boundary layer~\cite{monteith1965evaporation, idso1981normalizing}. Traditional precision irrigation methodologies rely on either static single-point soil moisture thresholding (e.g., triggering irrigation whenever volumetric soil water content drops below 20\%) or offline agro-hydrological water balance models~\cite{allen1998crop}. However, static thresholding inherently neglects the nonlinear thermodynamic interactions between edaphic moisture availability and atmospheric evaporative demand. On hot, arid afternoons characterized by steep vapor pressure deficits, a crop may experience severe cellular dehydration and turgor loss even when soil moisture appears superficially adequate. Conversely, on cool, humid, or overcast mornings, soil water depletion may remain benign without inducing transpirational stress.

To overcome the shortcomings of rule-based thresholds, recent research has turned to artificial intelligence and deep neural networks to infer water stress and environmental anomalies from multispectral, thermal, and sensor telemetry~\cite{jeaunita2025multitask, guo2025mffusionnet, choi2025sweetpotato}. However, state-of-the-art frameworks remain predominantly tethered to high-performance centralized cloud servers. For instance, the multi-task deep learning architecture proposed by Jeaunita \textit{et~al.}~\cite{jeaunita2025multitask} incorporates a 138-million-parameter VGG-16 backbone, InterPlanetary File System (IPFS) distributed storage, and multi-task deep autoencoders to achieve anomaly detection and emission forecasting in precision agriculture. While scientifically compelling, such cloud-dependent topologies face insurmountable bottlenecks when translated to actual commercial farming landscapes:
\begin{enumerate}
    \item \textbf{Telemetry and Connectivity Vulnerabilities:} Farmland and open orchards are frequently situated in rural cellular dead zones or regions afflicted by severe intermittent network dropouts. A cloud-tethered irrigation controller suffers catastrophic failure during communication disruptions, risking crop desiccation or perpetual valve flooding~\cite{olatinwo2025review}.
    \item \textbf{Energy and Bandwidth Overhead:} Continuous streaming of high-dimensional multi-spectral or spatial image streams over satellite or 4G/5G uplinks exhausts battery reserves in energy-harvesting field nodes within days and imposes unsustainable cellular data subscription costs~\cite{banbury2021benchmarking}.
    \item \textbf{Operational Latency vs. Actuation Real-Time Bounds:} Closed-loop valve actuation requires bounded, sub-second deterministic decision cycles that cannot tolerate the unpredictable packet delays and jitter inherent to remote server round-trips.
\end{enumerate}

To circumvent cloud dependency, the emerging discipline of Tiny Machine Learning (TinyML) seeks to compress machine intelligence directly onto edge microcontrollers (MCUs)~\cite{banbury2021benchmarking, lin2020mcunet, korostenskyi2025smart}. A notable edge-focused benchmark is AquaSense, introduced by Zaman \textit{et~al.}~\cite{zaman2026aquasense}, which deploys a LightGBM gradient boosted decision tree (GBDT) on an Espressif ESP32-S3 microcontroller coupled with LoRaWAN telemetry. While advancing on-device execution, AquaSense and similar edge attempts exhibit several critical limitations:
\begin{enumerate}
    \item \textbf{Prohibitive Sensor Costs:} AquaSense relies heavily on long-wave infrared (LWIR) radiometric thermal cameras (e.g., FLIR Lepton 3.5), significantly inflating deployment expenses and limiting widespread economic adoption across smallholder farms~\cite{nugroho2025lowcost, yang2026watermelon}.
    \item \textbf{Unengineered Feature Extraction:} Existing edge classifiers frequently ingest raw or sliding-window statistical features without embedding biophysical thermodynamic equations, necessitating deeper, wider decision trees that consume excessive Flash memory and SRAM~\cite{zaman2026aquasense}.
    \item \textbf{Class Imbalance and False-Negative Collapse:} Under real-world agricultural conditions, water deficit stress is a minority event relative to normal optimal conditions. Unweighted edge models suffer from severe recall collapse on drought stress, failing to detect moisture deficits until permanent wilting occurs.
    \item \textbf{Black-Box Opacity:} Microcontroller implementations lack local explainability (Tiny-XAI), preventing agronomists and autonomous controllers from verifying whether a diagnostic trigger stems from physical soil desiccation or transient sensor noise.
\end{enumerate}

To address these fundamental challenges, this paper presents \textbf{HydroEdge-Net}, a biophysically informed, explainable, ultra-lightweight TinyML edge computing framework for precision crop water stress diagnosis and closed-loop autonomous irrigation scheduling. HydroEdge-Net synthesizes fundamental soil-plant-atmosphere thermodynamic principles directly into edge feature preprocessing, replacing costly radiometric thermal imaging with accessible multi-parameter ambient and edaphic telemetry. The framework employs a cost-sensitively reweighted, depth-constrained gradient boosted decision tree architecture that compiles into bare-metal C code without dynamic memory allocation or interpreter dependencies.

The core contributions of this work are summarized as follows:
\begin{itemize}
    \item \textbf{Biophysically Informed Edge Feature Engine:} We formulate a lightweight on-device feature extraction pipeline that derives Tetens saturation vapor pressure $e_s(T)$, actual vapor pressure $e_a(T, \text{RH})$, atmospheric Vapor Pressure Deficit (VPD), Soil Moisture Deficit (SMD) relative to field capacity ($\theta_{\text{FC}} = 35.0\%$), and a novel Thermal-Moisture Stress Coupling Index ($\Psi_{\text{TMS}}$). By embedding nonlinear thermodynamic interactions and categorical encodings for seven soil textures and eight crop phenological stages, the system eliminates the need for expensive thermal cameras.
    \item \textbf{Cost-Sensitive TinyML Architecture:} We design a constrained LightGBM decision framework with strict leaf-wise depth limits ($d \le 5$, 31 leaves, 45 estimators) and dynamic class-weight compensation ($\omega = N_{\text{majority}} / N_{\text{minority}}$). On a benchmark of 16,411 multi-crop instances, HydroEdge-Net achieves 95.71\% overall accuracy and an exceptional 98.96\% sensitivity (recall) on critical water deficit stress.
    \item \textbf{Zero-Overhead Bare-Metal C Deployment:} We implement an automated compiler that translates trained decision trees into pure, statically linked C headers (\texttt{edge\_model.h}). The firmware requires zero heap allocation (\texttt{malloc}-free) and zero runtime dependencies (eliminating Treelite and TFLite-Micro overhead). On an Espressif ESP32-S3 (240~MHz), inference executes in 1.42~ms and consumes merely 16.8~KB of Flash and 11.4~KB of SRAM, representing a 40.0\% Flash reduction and 64.38\% SRAM reduction compared to the AquaSense baseline.
    \item \textbf{Tiny-XAI Physiological Explainability:} We incorporate TreeSHAP attribution analysis to explain feature importances and rigorously validate model behavior against physiological principles. The analysis confirms that soil moisture content (MOI), the proposed $\Psi_{\text{TMS}}$ index, VPD, and vulnerable early vegetative growth stages constitute the dominant diagnostic drivers, verifying biophysical consistency.
    \item \textbf{Closed-Loop Hydrological Digital Twin Validation:} We evaluate HydroEdge-Net in a 30-day (720-hour) continuous simulation under a dual-layer FAO-56 Penman-Monteith dynamic soil water balance model. HydroEdge-Net delivers a 63.96\% water conservation gain compared to calendar flooding while completely preventing transpirational water stress and soil hypoxia, whereas static thresholding under-irrigated and induced severe crop moisture stress.
\end{itemize}

The remainder of this paper is structured as follows. Section~\ref{sec:related_work} reviews related work and provides a gap analysis of cloud and edge architectures. Section~\ref{sec:methodology} introduces the biophysical mathematical formulation and the constrained GBDT model. Section~\ref{sec:experimental_setup} outlines the dataset, hardware testbed, and baseline models. Section~\ref{sec:results} presents benchmark performance, per-class evaluations, and confusion matrix comparisons. Section~\ref{sec:explainability} analyzes Tiny-XAI physiological explainability. Section~\ref{sec:digital_twin} details the 30-day closed-loop FAO-56 digital twin simulation. Section~\ref{sec:deployment} details on-device firmware compilation and silicon profiling. Section~\ref{sec:ablation} conducts a component ablation study. Finally, Section~\ref{sec:conclusion} concludes the paper and discusses future trajectories.
